\begin{proposition}\label{push:src:D.4.8}
For every \(H>H_*\), the degeneracy locus \(Z_H\) is exactly one
embedded \(C^1\) Jordan curve in the finite \(t\)-plane. It is
real analytic unless \(H=H_0=3\log48/(2\pi)\); at that height it
is real analytic away from \(t=0\) and has the regular crossing
described in \ref{push:src:D.4.5}. There are no singular coarse levels at any
supercritical height. The bounded complementary disc contains
\(t=1\) and has \(D_c>0\); the complementary disc in
\(\mathbb P^1\), containing the cusp, has \(D_c<0\). The point
\(t=0\) is in the latter disc for \(H<H_0\), on the curve for
\(H=H_0\), and in the former disc for \(H>H_0\).

The bounded discs are strictly nested as \(H\) increases, shrink
to \(\{1\}\) as \(H\downarrow H_*\), and exhaust the finite
plane as \(H\to\infty\). Their small-height boundary shape is
the expansion in \ref{push:src:D.4.1}. The maximal open base locus of the complex
manifold construction consists of these two complementary discs.
Approaching their common boundary, the covolume of the marked
root torus lattice is
\[
 \operatorname{covol}(\Pi_c\mathbb Z^4)
       =\operatorname{Im}\tau\,|D_c|\longrightarrow0.
\]
At the boundary the period map has rank three and the exact quotient
and Hausdorff reflection are those of \ref{push:src:D.4.3}. Thus this statement
classifies the locus and the local degeneration requested in D.4.
\end{proposition}
\Uses{\cref{ari:C.8.15}; \cref{push:src:D.4.1,push:src:D.4.2,push:src:D.4.3,push:src:D.4.5}; \cref{ari:C.1.1}}
\begin{proof}

Use the same exact affine relation and write
\[
 f=-D_{\rm part},\qquad
 Z_H=\{L=\log(2\sqrt3)-\pi H/6\}.
\]
The maximum $L(1)=\log\sqrt2$ corresponds exactly to
$H_*=3\log6/\pi$. Thus $H>H_*$ corresponds to a level
strictly below that maximum. Theorem~\ref{ari:C.8.15} proves
that each such level is one strictly star-shaped $C^1$ Jordan
curve about $1$, real analytic away from the possible point
$t=0$. Its nonzero coarse derivative there proves regularity
also on that exceptional level.

Since $L(0)=\tfrac14\log3$, its height is
$H_0=3\log48/(2\pi)$. The exact equation $D_c=H-f$ gives
the two signs, the position of $0$, and the bounded component
$\{f<H\}$ containing $1$. Strict radial monotonicity proves
strict nesting. The unique minimum $f(1)=H_*$ and properness
of $f$ off $1$ prove shrinking as $H\downarrow H_*$ and
exhaustion of every compact subset as $H\to\infty$.
The expansion at the minimum is \ref{push:src:D.4.1}.

In the ordered real coordinates and marked lattice basis of
\ref{push:src:D.4.3}, direct expansion of the period determinant
gives
\[
 \det_{\R}\Pi_c
   =\Im\tau\,\Im\beta_c-6(\Im\mu)^2
   =\Im\tau\,D_c.
\]
Its absolute value is the covolume and tends to zero at the
level curve. The same matrix retains three independent columns
there, so its rank is exactly three. Apply the exact quotient
and Hausdorff-reflection computation of \ref{push:src:D.4.3}.
The local proper-action criterion of \ref{push:src:D.4.2}
gives complex manifolds on both complementary regions and
excludes every point of their common boundary, including its
discrete fixed-fibre quotients. This proves the maximality and
all the local degeneration assertions.

\end{proof}
